%% ===========================================================================
\FloatBarrier
\section{Benchmarks and convergence study}
\label{app:verification}
\setcounter{table}{0}
\setcounter{figure}{0}
%% ===========================================================================

This appendix collects the benchmarks behind Sec.~\ref{sec:verification} in
full: the benchmark, identity and independent-route record, and the
convergence study. The numerical precision of the thinnest point follows the
two tables.

Table~\ref{tab:verification} lists the benchmark problems, the exact
identities and the tolerance each was required to meet, together with the value
obtained. The identities are the ones that a re-implementation can reproduce
without any of the data of this paper: the patch-test energy of a constant
strain field, the second-moment coefficients of \ref{app:Mcoeff}
recovered from a numerical integration of the window, the closed-form coupling
factor of a uniaxial crystal, and the analytical damping ratio of the classical
Kelvin--Voigt element. The rows marked as independent routes are values
recomputed along a path that does not use the assembled matrices at all --
Euler--Bernoulli theory for the bar frequency, the constitutive tensors for the
coupling factor, the closed forms of \ref{app:Mcoeff} for the window
coefficients, and the analytical limit for the damping ratio.

\begin{table}[!htbp]
\centering
\small
\caption{Verification record: benchmarks, exact identities and independent
routes, each with the value obtained and the tolerance it was required to
meet. The material, the geometry and the mesh of each row are those of
Sec.~\ref{sec:material}.}
\label{tab:verification}
%% A compact list of the numerical benchmarks: what was computed, what it is
%% compared with, and the agreement.
\begin{tabular}{@{}l l l@{}}
\toprule
check & obtained & reference \\
\midrule
patch test, relative energy error & $4.44\times10^{-16}$ & $0$ \\
gradient energy at constant strain & $1.06\times10^{-12}$ & $0$ \\
$K_g$ of the ellipsoid vs closed form & $3.40\times10^{-16}$ & $0$ \\
$K_g$(cuboid)$/K_g$(ellipsoid), equal half-width & 1.666667 & $5/3$ \\
equal-$M$ cuboid and ellipsoid, max rel.\ difference & $0$ & $0$ \\
tip displacement, length-scale sweep & 4.2339\,\si{\micro\meter} & 4.23\,\si{\micro\meter} \\
observed order of convergence & 1.68 & $2.01$ on the finest pair \\
piezoelectric $k^2$ of extensional mode 5, plane strain vs plane
stress & 0.4024 vs 0.4580 & mode 5 of each model, MAC $0.99999$ \\
L1 order at $\alpha=1$ & 0.99 & $1$ \\
free strip $k^2$, FEM vs 2-D closed form & 0.56079 vs 0.56079 & $<10^{-3}$ rel. \\
free strip, 3-D vs 2-D reduction & 2.83\,\% & modeling gap \\
clamped strip, $96\times4$ vs $48\times4$ & 0.55992 vs 0.55957 & $6.2\times10^{-4}$ rel. \\
flexoelectric energy, uniform curvature & $3.570415\times10^{-4}$ & $3.570415\times10^{-4}$ \\
open-circuit voltage, uniform curvature & $1.7\times10^{-13}$\,\si{\volt} & $<10^{-10}$\,\si{\volt} \\
clamped-free $f_1$, FEM vs Euler--Bernoulli & 144.6 vs 143.8\,\si{\kilo\hertz} & $128\times4$: 144.08 \\
clamped-clamped $f_1$, FEM vs Euler--Bernoulli & 919.2 vs 915.0\,\si{\kilo\hertz} & $48\times4$ \\
modal truncation, 24 modes & $3.19\times10^{-2}$ & monotone in $N$ \\
\bottomrule
\end{tabular}

\end{table}

The two closed forms of Sec.~\ref{sec:verification} rest on the scalar
$e^{\mathsf T}C^{-1}e=6.3308\times10^{-10}$ of the reduced tensors, and the
\SI{0.5}{\percent} agreement between them and the finite element value is the
agreement of that scalar with the assembled route rather than a coincidence of
two discretizations.

The clamped--free frequency is the one row with an exact counterpart, so it is
also reported under refinement: the $128\times4$ mesh gives
144.08\,\si{\kilo\hertz}, against 144.6\,\si{\kilo\hertz} on the
$48\times4$ mesh and 143.8\,\si{\kilo\hertz} from Euler--Bernoulli theory.

The two-dimensional reduction is quantified in the same record: the
plane-stress coupling differs from the three-dimensional value by $2.8\,\%$,
which is the size of the modeling error that the reduction introduces and the
reason the two reductions are not used interchangeably. Of the independent
values, the damping ratio of the element at $\alpha=0.3$ and $\omega\tau=0.3$
is $0.108$, against a peak of $0.179$ at $\alpha=0.7$
(Sec.~\ref{sec:memorystudy}).

Two rows of Table~\ref{tab:verification} are comparisons rather than exact
identities. The order of convergence is the least-squares slope of
$\log|u_R/u_{\rm ref}-1|$ against $\log h$, with $h=L/n_x$ halved over the
four meshes $5\times1$, $10\times2$, $20\times4$ and $40\times8$ of the
gradient-free benchmark of Sec.~\ref{sec:verification}; the slope over all
four is 1.68 and the finest pair alone gives 2.01, so the sequence is entering
its asymptotic range from the coarse end. The expected rate for the bicubic
Hermite rectangle is $O(h^{4})$ in $L^2$ and $O(h^{3})$ in $H^1$ on a smooth
solution, and a linear functional such as $u_R$ converges at least as fast as
the energy norm; the requirement is the empirical slope above 1.5, and no
theoretical exponent is claimed beyond that bound. The extensional coupling is
compared across the two reductions on the same $24\times2$ mesh: $0.4024$ in
plane strain against $0.4580$ in plane stress. The two models are compared
mode by mode before the couplings are, because the two reductions could in
principle select different shapes and then the comparison would be empty. They
do not: the shape selected by the rule of Table~\ref{tab:application} is mode
5 in both, and the two short-circuit mode shapes, mass-normalized, overlap
with $\mathrm{MAC}=0.99999$, while the largest overlap between different
members of the two six-mode sets is $2.8\times10^{-6}$. The coupling therefore
differs by the elastic reduction alone, which stiffens the section and lowers
$k^{2}=e^{2}/(C\kappa)$; the difference is a property of the model.

Table~\ref{tab:convergence} is the convergence study behind
Fig.~\ref{fig:convergence}. The frequency of the extensional mode and its
mode-free coupling factor are listed for four meshes, with the errors referred
to the $48\times4$ mesh used in the paper: the coarser $24\times2$ mesh of the
parameter sweeps differs from it by $0.04\,\%$ in frequency and by
$0.03\,\%$ in the coupling factor, and the other two rows by less. None of
the four rows is a clamped strip; the
clamped strip is the $96\times4$ against $48\times4$ entry of
Table~\ref{tab:verification}, whose coupling factors differ by
$6.2\times10^{-4}$ relative.

%% generated table -- do not edit (see AUTHOR_TODO.md)
\begin{table}[!htbp]
\centering
\caption{Convergence of the extensional mode behind
Fig.~\ref{fig:convergence}: strip \SI{200}{\micro\meter}$\times$\SI{10}{\micro\meter},
poling $x$, plane stress, $\mu=0$, $a_1\times a_2=3\times1$\,\si{\micro\meter}.
The frequency is that of the extensional mode of the short-circuit eigenvalue
problem and the coupling is the mode-free quasi-static value of
\eqref{eq:kstatic}, the number quoted in the text; the errors are referred to
the $48\times4$ mesh.}
\label{tab:convergence}
\begin{tabular}{@{}l r r r r r@{}}
\toprule
mesh & dof & $f$ (kHz) & $k^2$ & err $f$ & err $k^2$ \\
\midrule
$12\times2$ & 312 & 4745.4 & 0.55682 & 3.2e-03 & 2.4e-03 \\
$24\times2$ & 600 & 4732.0 & 0.55797 & 3.7e-04 & 3.0e-04 \\
$36\times2$ & 888 & 4730.0 & 0.55814 & 3.7e-05 & 7.0e-06 \\
$48\times4$ & 1960 & 4730.2 & 0.55814 & 0.0e+00 & 0.0e+00 \\
\bottomrule
\end{tabular}
\end{table}


\paragraph{Numerical precision of the thinnest point} The
$h=\SI{1}{\micro\meter}$ coupling is a difference of two nearly equal
eigenvalues of a pencil whose condition number is $\sim10^{12}$, so it
inherits the accuracy of the eigensolver. The same model was solved through
four mathematically equivalent eigensolver paths -- a QZ generalized
eigenproblem, a dense symmetric reduction, a sparse Lanczos reduction and a
truncated modal sum -- and through four numeric stacks, and all of them return
$4.2\times10^{-3}$ and the factor $24$ to within $0.5\,\%$. Deviations are
measured from the midpoint of the equivalent results, so the spread quoted is
the spread of the model, and at $h=\SI{10}{\micro\meter}$ it is at most
$10^{-4}\,\%$. Every path keeps the whole spectrum in the secular sum:
truncating it to twelve modes moves $k^{2}$ by more than $80\,\%$, an effect
of the approximation alone. A run of the same suite on an additional numeric stack
returns $4.18\times10^{-3}$ ($+0.2\,\%$), consistent with this statement. This
spread fixes the two significant digits quoted in the text and in the
abstract.

\FloatBarrier
